\chapter{Машинне навчання та статистичне виявлення}
\chaptershort{Машинне навчання та статистичне виявлення}
\label{chap:ch09}

Експоненційне зростання кількості зразків (розділ~\ref{chap:ch01}) та вичерпність суто сигнатурного підходу перед поліморфізмом і метаморфізмом (розділ~\ref{chap:ch08}) зробили \term{машинне навчання} (МН) центральним інструментом сучасного виявлення ЗПК. Замість ручного опису ознак кожної загрози МН \emph{навчається} розрізняти злоякісне й легітимне за прикладами, узагальнюючи на небачені зразки. Це дає виявлення нових загроз (на що нездатні сигнатури) і масштабованість на мільйони зразків. Проте МН не є чарівним рішенням: воно успадковує й загострює проблему базової частоти (розділ~\ref{chap:ch03}), потребує якісних розмічених даних, вразливе до змагальних атак (розділ~\ref{chap:ch08}) і дає рішення, які важко пояснити. Розділ розглядає місце МН у виявленні, ознаки й подання зразків, типові моделі, оцінювання якості, змагальне машинне навчання та практичні обмеження, і завершується кейсом побудови статичного класифікатора PE-файлів. МН не замінює інших методів, а інтегрується з ними (розділи~\ref{chap:ch06}, \ref{chap:ch07}, \ref{chap:ch10}).

\section{Місце машинного навчання у виявленні}\label{sec:ch09-place}

МН у протидії ЗПК розв'язує дві основні задачі: \emph{класифікацію} (злоякісний/легітимний --- бінарна; або віднесення до сімейства --- багатокласова) і \emph{кластеризацію} (виявлення груп схожих зразків без розмітки, розділ~\ref{chap:ch06}). За наявністю розмітки розрізняють навчання з учителем (на розмічених прикладах), без учителя (кластеризація, виявлення аномалій) та напівкероване. Фундаментальна ідея --- узагальнення: модель, навчена на відомих зразках, має правильно класифікувати \emph{нові}, яких не було в навчальній вибірці. Саме здатність до узагальнення відрізняє МН від сигнатур і дає виявлення небачених загроз~\cite{ch09:ucci,ch09:gibert}.

Класична праця~\cite{ch09:sommer} застерігає від наївного перенесення МН із задач на кшталт розпізнавання образів у безпеку: тут висока ціна хибних тривог, величезний дисбаланс класів, змінність (concept drift --- загрози еволюціонують, тож модель старіє), активний супротивник, що цілеспрямовано обходить модель. Ці особливості визначають, чому МН у безпеці складніше, ніж у типових застосуваннях, і чому воно доповнює, а не замінює детерміновані методи.

Корисно розрізняти два режими застосування МН, що мають різні вимоги й ризики. \emph{Хмарне} виявлення на боці постачальника працює з величезними обчислювальними ресурсами та агрегованою телеметрією з мільйонів вузлів; воно може дозволити собі дорогі моделі й швидке перенавчання, але вносить затримку й залежність від зв'язку. \emph{Локальне} виявлення на кінцевій точці (розділ~\ref{chap:ch10}) працює в реальному часі з обмеженими ресурсами й без гарантованого зв'язку; воно потребує компактних, швидких моделей і особливо чутливе до хибних тривог, бо помилка блокує роботу користувача негайно. Промислові рішення поєднують обидва режими: легка локальна модель дає миттєве рішення, а підозрілі чи невизначені випадки ескалюються в хмару для глибшого аналізу. Таке розшарування --- ще один прояв каскадної архітектури із зростанням вартості й покриття, що проходить крізь увесь посібник.

Важливо усвідомлювати межу застосовності узагальнення. Модель надійно впізнає варіації, \emph{подібні} до баченого в навчанні, але принципово нова техніка, не представлена у вибірці, може лишитися невиявленою так само, як і сигнатурою. МН не скасовує гонитви озброєнь, а зміщує її: замість створення нової сигнатури під кожен зразок захисник підтримує актуальність навчальної вибірки й моделі під еволюцію цілих класів загроз. Виграш --- у масштабі узагальнення, а не в принциповій повноті, яка, за результатами розділу~\ref{chap:ch03}, недосяжна для жодного методу.

\section{Ознаки та подання зразків}\label{sec:ch09-features}

Якість МН визначається насамперед \term{ознаками} --- числовим поданням зразка, що подається моделі. Ознаки поділяють на статичні (з файлу без виконання) і динамічні (з поведінки, розділ~\ref{chap:ch07}). Типові статичні ознаки: байтові $n$-грами, частоти інструкцій, імпортовані функції, властивості PE-заголовка й секцій, ентропія~\eqref{eq:ch03-entropy}, рядки, розмір. Піонерська праця~\cite{ch09:schultz} уперше застосувала здобуття даних до цих ознак; відкритий набір EMBER~\cite{ch09:ember} стандартизував ознаки PE для досліджень. Динамічні ознаки: послідовності викликів API, дії над ресурсами, мережева активність.

\begin{tikzfig}
\begin{tikzpicture}[font=\scriptsize,
  b/.style={draw,rounded corners=2pt,minimum height=6mm,text width=20mm,align=center,font=\tiny},
  >={Stealth[length=1.5mm]}]
\node[b] (s) at (0,0) {зразок};
\node[b] (f) at (2.6,0) {виділення ознак (вектор)};
\node[b] (m) at (5.2,0) {модель};
\node[b] (o) at (7.6,0) {ЗПК / чисте + певність};
\draw[->] (s)--(f); \draw[->] (f)--(m); \draw[->] (m)--(o);
\node[b,fill=black!6] (tr) at (5.2,-1.3) {навчальна вибірка};
\draw[->] (tr)--(m);
\end{tikzpicture}
\figcap{Рис.~\thechapter.1 --- Конвеєр виявлення на основі машинного навчання: від зразка через виділення ознак до рішення моделі}
\end{tikzfig}

Окремий напрям --- навчання \emph{без} ручного виділення ознак: глибокі моделі приймають «сирі» байти~\cite{ch09:raff} або двовимірне подання файлу~\cite{ch09:saxe} й самі навчаються ознак. Це знижує залежність від експертних знань, але потребує більших даних і обчислень та дає менш інтерпретовані моделі. Вибір між ручними й автоматичними ознаками --- компроміс між інтерпретовністю, вартістю й стійкістю до ухилення.

Сирі ознаки зазвичай потребують \emph{нормалізації} та зважування, бо різні ознаки мають різні масштаби й інформативність. Класичний прийом для текстоподібних ознак (рядків, $n$-грамів інструкцій) --- зважування за частотою появи в зразку відносно поширеності в усій колекції (TF--IDF), що підкреслює рідкісні, а отже дискримінативні, ознаки й пригнічує всюдисущі. Для ознаки $t$ у зразку $d$ з колекції $D$ вага

\begin{equation}
w_{t,d} = f_{t,d}\cdot\log\frac{|D|}{|\{d'\in D: t\in d'\}|},
\label{eq:ch09-tfidf}
\end{equation}

де $f_{t,d}$ --- частота ознаки в зразку, а логарифмічний множник (зворотна документна частота) тим більший, чим рідкісніша ознака. Формула~\eqref{eq:ch09-tfidf} ілюструє загальний принцип інженерії ознак: цінність ознаки визначає не частота сама по собі, а здатність \emph{розрізняти} класи. Погано дібрані або надлишкові ознаки не лише не допомагають, а й підвищують ризик перенавчання та відкривають простір для змагального ухилення.

Окрему увагу приділяють \emph{витоку} інформації між навчанням і тестом. Поширена методологічна помилка --- включення до навчальної вибірки варіантів тих самих зразків, що й у тесті, або ознак, які опосередковано кодують мітку (наприклад, артефактів конкретного аналітичного конвеєра). Такий витік дає оптимістично завищену оцінку якості, що не відтворюється в експлуатації. Тому коректне формування вибірок --- не технічна дрібниця, а передумова вірогідності всіх подальших висновків.

\section{Типові моделі та їх властивості}\label{sec:ch09-models}

У виявленні ЗПК застосовують широкий спектр моделей, що різняться інтерпретовністю, вартістю й стійкістю (табл.~\thechapter.1). \emph{Дерева рішень та їх ансамблі} (випадковий ліс, градієнтний бустинг) --- робочий інструмент статичного виявлення PE: вони добре працюють із табличними ознаками, відносно інтерпретовані й стійкі. \emph{Лінійні моделі} (логістична регресія, SVM) прості й швидкі, але обмежені в складності меж. \emph{Нейронні мережі}, зокрема глибокі, здатні навчатися складних залежностей і працювати із сирими даними, але потребують великих вибірок, обчислень і важко інтерпретуються. \emph{Методи виявлення аномалій} шукають відхилення від норми, корисні для невідомих загроз, але схильні до хибних тривог.

\tabhead{Таблиця~\thechapter.1 --- Моделі машинного навчання у виявленні ЗПК}

\begin{center}\footnotesize
\begin{tabularx}{\linewidth}{L{24mm} Y L{18mm}}
\toprule
Модель & Сильні риси & Інтерпретовність \\
\midrule
Ансамблі дерев & точність на табличних ознаках & помірна \\
Лінійні (SVM) & швидкість, простота & висока \\
Нейромережі & складні залежності, сирі дані & низька \\
Виявлення аномалій & невідомі загрози & помірна \\
\bottomrule
\end{tabularx}
\end{center}

Немає універсально найкращої моделі: вибір залежить від задачі, доступних даних, вимог до швидкості та інтерпретовності й моделі загроз. На практиці ефективні ансамблі та багаторівневі архітектури, що поєднують кілька моделей і типів ознак, а також інтеграція МН у каскад із сигнатурами й поведінковим аналізом.

Особливу увагу в безпеці приділяють \emph{інтерпретовності} моделі. Коли детектор блокує файл, аналітикові центру реагування (розділ~\ref{chap:ch16}) потрібно зрозуміти \emph{чому}, щоб підтвердити рішення, розслідувати інцидент і відрізнити справжнє спрацювання від хибного. Інтерпретовані моделі (дерева, лінійні) дають таке пояснення безпосередньо через ваги ознак чи шлях у дереві; для складних моделей застосовують методи пояснення post-hoc, що наближено вказують, які ознаки найбільше вплинули на рішення. Брак інтерпретовності --- не лише операційна незручність, а й регуляторний ризик (розділ~\ref{chap:ch02}) та перешкода довірі, тому в критичних застосуваннях часто свідомо жертвують часткою точності заради прозорості рішень.

\section{Оцінювання якості}\label{sec:ch09-eval}

Оцінювання детектора на основі МН спирається на метрики розділу~\ref{chap:ch03}: повноту, точність, частоту хибних тривог, площу під ROC. Критична особливість --- величезний дисбаланс і ціна хибних тривог: за мільйонів легітимних файлів навіть мала частка хибних тривог дає неприйнятний абсолютний потік (проблема базової частоти). Тому в безпеці оцінюють якість за дуже низьких рівнів хибних тривог. Для бінарного класифікатора з розмітками «злоякісний» ($+$) і «чистий» ($-$) зведену міру якості дає $F_\beta$-міра

\begin{equation}
F_\beta = (1+\beta^2)\cdot\frac{\mathrm{Precision}\cdot\mathrm{Recall}}{\beta^2\cdot\mathrm{Precision}+\mathrm{Recall}},
\label{eq:ch09-fbeta}
\end{equation}

де параметр $\beta$ задає відносну вагу повноти й точності: $\beta<1$ наголошує на точності (мінімізації хибних тривог), що відповідає пріоритетам захисту кінцевих точок. Коректне оцінювання вимагає реалістичного складу тестової вибірки (базова частота як у продакшні), часового поділу (навчання на минулому, тест на майбутньому --- щоб урахувати дрейф) і незалежності тесту від навчання.

Окрема небезпека --- \emph{дрейф концепту}: загрози еволюціонують, тому модель, точна сьогодні, деградує з часом. Це вимагає моніторингу якості в експлуатації та періодичного перенавчання на свіжих даних, а також оцінювання не лише миттєвої, а й \emph{часової} стійкості моделі.

Практичним наслідком пріоритету низьких хибних тривог є вибір \emph{робочої точки} на ROC-кривій (розділ~\ref{chap:ch03}). Модель видає не бінарне рішення, а оцінку певності; поріг, за яким оцінку перетворюють на вердикт, визначає компроміс між повнотою й хибними тривогами. У захисті кінцевих точок поріг зсувають так, щоб частота хибних тривог не перевищувала жорсткого операційного ліміту (наприклад, одиниці на мільйони файлів), свідомо жертвуючи частиною повноти. Це означає, що «точність моделі» без указання робочої точки й базової частоти тестового набору є беззмістовною характеристикою: та сама модель за одного порогу придатна до промислового розгортання, а за іншого --- генерує неприйнятний потік хибних тривог. Тому в безпеці звітують не одне число, а поведінку детектора в діапазоні низьких хибних тривог, і саме там порівнюють конкурентні моделі.

\section{Змагальне машинне навчання}\label{sec:ch09-adversarial}

Протидія ЗПК --- змагальне середовище з активним супротивником, тому МН тут підлягає \term{змагальним атакам}~\cite{ch09:biggio}. \emph{Ухилення} (evasion): зловмисник мінімально змінює зразок так, щоб модель класифікувала його як чистий, зберігаючи зловмисну функцію~\cite{ch09:grosse}. На відміну від задач розпізнавання образів, тут зміни обмежені вимогою \emph{працездатності} коду, тому застосовують додавання «нешкідливих» секцій, зміну метаданих, вставлення легітимних ознак. Показано, що навіть навчанням із підкріпленням можна автоматично знаходити модифікації, що обходять статичні моделі PE~\cite{ch09:anderson-gym}. \emph{Отруєння} (poisoning): зловмисник впливає на навчальні дані, щоб спотворити модель.

\begin{tikzfig}
\begin{tikzpicture}[font=\scriptsize,
  b/.style={draw,rounded corners=2pt,minimum height=6mm,text width=22mm,align=center,font=\tiny},
  >={Stealth[length=1.5mm]}]
\node[b] (m) at (0,0) {зразок ЗПК};
\node[b] (p) at (3,0) {+ нешкідливі зміни (зберігають функцію)};
\node[b] (mod) at (6,0) {модель: «чистий»};
\draw[->] (m)--(p); \draw[->] (p)-- node[above,font=\tiny]{ухилення} (mod);
\node[b,fill=black!6] (d) at (3,-1.3) {контрзаходи: змагальне навчання, ансамблі, межі};
\draw[->,dashed] (d)--(mod);
\end{tikzpicture}
\figcap{Рис.~\thechapter.2 --- Змагальне ухилення від класифікатора й напрями протидії}
\end{tikzfig}

Контрзаходи: змагальне навчання (додавання змагальних прикладів до навчальної вибірки), ансамблі різнорідних моделей і ознак (ухилення від усіх одразу складніше), опора на ознаки, які важко змінити без втрати функції (семантичні, поведінкові), моніторинг аномалій у вхідних даних. Фундаментальний висновок: у змагальному середовищі не буває остаточно «навченої» моделі --- потрібен безперервний цикл оцінювання, оновлення й диверсифікації. Це перегукується зі стратегічним принципом розділу~\ref{chap:ch08}: виявлення слід будувати на ознаках, інваріантних до дешевих модифікацій.

Корисно усвідомлювати асиметрію змагальної задачі. Зловмисникові достатньо знайти \emph{одну} модифікацію, що обходить модель, тоді як захисник мусить встояти проти \emph{всіх} можливих. Ця асиметрія --- та сама, що й у безпеці загалом, --- означає, що абсолютна стійкість моделі недосяжна, а метою є підвищення вартості ухилення до рівня, за якого воно стає невигідним. Практично це досягається поєднанням: по-перше, опорою на ознаки, зміна яких руйнує функцію ЗПК (семантичні, поведінкові), бо тут простір дешевих модифікацій вужчий; по-друге, ансамблюванням різнорідних моделей, що змушує зловмисника обходити кілька незалежних детекторів одночасно; по-третє, таємністю деталей моделі (хоча на неї не варто покладатися як на єдиний захист, бо модель можна досліджувати запитами). Варто також розрізняти ухилення в \emph{білому} ящику (зловмисник знає модель і може обчислювати градієнти) та в \emph{чорному} (лише спостерігає вердикти): перше дає потужніші атаки, друге реалістичніше для більшості сценаріїв і потребує багатьох запитів, що саме може бути виявлене як аномалія.

Отруєння навчальних даних --- загроза, специфічна для конвеєрів, що автоматично поповнюють вибірку з відкритих джерел (наприклад, краудсорсингових колекцій зразків). Зловмисник, здатний впливати на ці джерела, може підсунути хибно розмічені або спеціально сконструйовані зразки, що зміщують межу рішення на його користь --- створюючи «сліпу пляму» під конкретне сімейство. Контрзаходи --- контроль якості й походження даних, стійкі до викидів методи навчання, моніторинг раптових змін у розподілі навчальної вибірки. Отруєння ілюструє, що в безпеці довіряти не можна навіть власним навчальним даним, якщо супротивник здатен на них впливати.

\section{Практичні обмеження й інтеграція}\label{sec:ch09-limits}

Практичне застосування МН обмежують кілька чинників. \emph{Дані}: потрібні великі, якісні, актуальні, коректно розмічені вибірки; хибна розмітка й зміщення вибірки псують модель. \emph{Інтерпретовність}: рішення складних моделей важко пояснити, що ускладнює довіру, розслідування й відповідність регуляторним вимогам (розділ~\ref{chap:ch02}). \emph{Дрейф і супровід}: моделі старіють і потребують перенавчання. \emph{Змагальність}: активний супротивник. \emph{Ресурси}: навчання й інференс глибоких моделей дорогі. Через це МН у промислових засобах не застосовують ізольовано, а інтегрують у багаторівневий конвеєр із сигнатурами, репутацією та поведінковим аналізом (розділи~\ref{chap:ch06}, \ref{chap:ch07}, \ref{chap:ch10}), де кожен рівень компенсує слабкості інших. Вітчизняні дослідження застосування МН~\cite{ch09:ml-ua} та нейромережевої класифікації за статичними ознаками~\cite{ch09:neural-ua} підтверджують ці висновки на українських даних.

\section{Кейс: статичний класифікатор PE-файлів}\label{sec:ch09-case}

Розглянемо побудову статичного класифікатора «злоякісний/чистий» для PE-файлів (захисна задача).

\emph{Крок 1. Дані.} Формуємо збалансований за часом набір розмічених зразків (злоякісних і легітимних) з реалістичним складом, із часовим поділом на навчання й тест.

\emph{Крок 2. Ознаки.} Виділяємо статичні ознаки: властивості заголовка й секцій, ентропію~\eqref{eq:ch03-entropy}, імпорти, байтові гістограми, рядки (за зразком EMBER~\cite{ch09:ember}).

\emph{Крок 3. Модель.} Навчаємо ансамбль дерев (градієнтний бустинг) як компроміс точності, швидкості й інтерпретовності.

\emph{Крок 4. Оцінювання.} Вимірюємо повноту за фіксованого низького рівня хибних тривог, ROC та $F_\beta$~\eqref{eq:ch09-fbeta} з $\beta<1$. Оцінюємо часову стійкість на пізніших зразках (дрейф).

\emph{Крок 5. Стійкість і супровід.} Перевіряємо стійкість до змагального ухилення (додавання нешкідливих ознак), за потреби застосовуємо змагальне навчання й ансамблювання. Плануємо моніторинг якості й перенавчання.

Кейс демонструє, що ефективний класифікатор --- не лише вдала модель, а дисциплінований процес: реалістичні дані, коректне оцінювання за низьких хибних тривог, урахування дрейфу й змагальності, інтеграція в загальну архітектуру захисту.

\section*{Висновки до розділу}
\addcontentsline{toc}{section}{Висновки до розділу}

Машинне навчання стало центральним інструментом виявлення ЗПК, бо дає узагальнення на нові загрози й масштабованість, яких бракує сигнатурам. Воно розв'язує задачі класифікації й кластеризації, спираючись на статичні та динамічні ознаки, причому якість визначається насамперед ознаками, а не лише моделлю. Спектр моделей --- від інтерпретованих ансамблів дерев до глибоких мереж на сирих даних --- не має універсального лідера; вибір залежить від задачі, даних і вимог (табл.~\thechapter.1). Оцінювання в безпеці підпорядковане проблемі базової частоти: якість міряють за дуже низьких рівнів хибних тривог~\eqref{eq:ch09-fbeta}, з урахуванням дрейфу концепту й часової стійкості. Протидія ЗПК --- змагальне середовище, тому МН вразливе до ухилення й отруєння, що вимагає змагального навчання, ансамблів, опори на стійкі ознаки й безперервного циклу оновлення. Практичні обмеження --- дані, інтерпретовність, дрейф, змагальність, ресурси --- роблять ізольоване застосування МН недоцільним. Головний висновок: МН --- потужний, але не самодостатній компонент; його цінність розкривається в інтеграції з сигнатурами, репутацією та поведінковим аналізом у багаторівневому конвеєрі з дисциплінованим супроводом.

\selfcheck
\begin{enumerate}
\item Чому машинне навчання стало центральним інструментом виявлення ЗПК? Яку властивість сигнатур воно долає?
\item Які задачі розв'язує МН у протидії ЗПК? Поясніть ідею узагальнення.
\item Від чого застерігає праця~\cite{ch09:sommer}? Назвіть особливості безпеки, що ускладнюють застосування МН.
\item Що таке ознаки? Порівняйте статичні й динамічні ознаки та ручне й автоматичне їх виділення.
\item Опишіть типові моделі (табл.~\thechapter.1) та їх компроміси. Чому немає універсально найкращої моделі?
\item Чому в безпеці якість оцінюють за дуже низьких рівнів хибних тривог? Як це пов'язано з базовою частотою (розділ~\ref{chap:ch03})?
\item Поясніть $F_\beta$-міру~\eqref{eq:ch09-fbeta}. Чому для захисту кінцевих точок обирають $\beta<1$?
\item Що таке дрейф концепту? Як він впливає на оцінювання й супровід моделі?
\item Що таке змагальне ухилення? Чим воно обмежене у випадку ЗПК порівняно з розпізнаванням образів?
\item Що таке отруєння навчальних даних? Назвіть контрзаходи проти змагальних атак.
\item Чому в змагальному середовищі не буває остаточно «навченої» моделі?
\item Які практичні чинники обмежують застосування МН? Чому МН інтегрують у багаторівневий конвеєр?
\item Опишіть побудову статичного класифікатора PE (кейс підрозділу~\ref{sec:ch09-case}). Чому важливий часовий поділ даних?
\item Дослідницьке: оцініть деградацію якості обраної моделі в часі (дрейф) на наборі зразків, розподілених за датами, і запропонуйте стратегію перенавчання.
\item Дослідницьке: порівняйте стійкість моделей на ручних та на сирих ознаках до автоматичного змагального ухилення.
\end{enumerate}
